%!TEX root = ../thesis.tex

\chapter{Experimental Methodology}\label{chap:methodology}

Due to the high variance of RL algorithms, we run each experiment using four different seeds.
We first show that the input-action targets can be successfully learned.
This serves as a correctness test for our EfficientZero implementation.
Since we will mainly focus on corpus-level actions, we will not do a parameter study,
but instead stop once the target was successfully learned.
For the corpus-level action targets, we will do a detailed ablation study on the cjson target.
We then use its parameters to test two additional targets, libxml2 as well as picohttpparser.

Experiment results are extracted using hooks (see~\cref{hooks}), either
during training for metrics like losses, or after training and during
evaluation for more expensive data collection like seed persistence.
Because of our architecture, we can easily evaluate new targets by using the existing
hook and evaluation infrastructure.

% Everything about HOW results are produced and judged lives here, so that
% \cref{chap:results} can be pure findings. Written before the results are read, and
% not revised to fit them --- say so.

\section{Experimental Setup}\label{sec:methodology-setup}

The Burn deep learning framework itself supports multiple backends.
We run our experiments on devices that all have access to NVIDIA GPUs and therefore use
the CUDA backend.
Experiments were run on two VMs with an Intel(R) Xeon(R) Gold 5418Y processor as
well as an NVIDIA L40S and 500 GiB of RAM.
This allows to run about two to four experiments of a campaign at once.

\subsection{Parameters}

There are many possible parameters to tune in this work.
We decide to mostly adopt the default EfficientZero parameters without any further studies
since the authors have already done a detailed investigation and found these to be good.
Two things we change is the temperature schedule $T:\mathbb{N}\rightarrow\mathbb{R}$ as well as
the root Dirichlet noise of \cref{eq:ez-dirichlet}, both its concentration $\xi$ and the fraction
$\rho$ it is mixed into the priors with, since we found fuzzing with large action spaces to require
more exploration.
The default EfficientZero parameters we adopt can be found in \cref{tab:app-ez-params}.
The remaining differences to the reference configuration are consequences of the scale of our
targets rather than of tuning: our observations are flat coverage vectors instead of stacked Atari
frames, and our rewards live on a much smaller scale (\cref{sec:mdp-rewards}).
\Cref{sec:methodology-hyperparameters} justifies each of them.

% TODO: The run configuration that \cref{chap:implementation} deliberately left out:
% hardware, backend, software versions, what a single run costs in wall-clock, and
% how campaigns are launched and terminated. Enough for someone else to repeat it.

\section{Metrics}\label{sec:methodology-metrics}

% TODO: Define every metric before it appears in a plot, grouped by what it is for.
%
% Headline metrics:
%   - maximum prefix depth reached per episode, over environment steps --- the most
%     legible learning curve on the diagnostic targets,
%   - solve rate over training, from a greedy noise-free evaluation rollout rather
%     than inferred from noisy training returns,
%   - first-passage time per depth and executions to first and to stable solve ---
%     the sample-efficiency headline,
%   - coverage versus executions, the standard fuzzing axis, for every approach on
%     the same plot,
%   - return and episode length.
%
% Mechanism metrics (these answer *why*, and are what the probe suite of
% \cref{sec:impl-probes} exists for):
%   - per-depth value and value-prefix ranking of the correct action,
%   - policy prior versus depth over training, and the policy-collapse matrix,
%   - MCTS visit-count entropy per depth,
%   - policy improvement: divergence between the MCTS visit-count policy and the
%     network prior at the root --- is search improving on the prior or echoing it?
%   - value calibration: predicted root value against realized return.

\section{Baselines}\label{sec:methodology-baselines}

% TODO: Three reference points, always plotted on the same axes:
%   - uniform random action selection: the floor,
%   - a coverage-guided greedy climber that keeps coverage-increasing prefixes but has
%     no learned model: the CONTROL that isolates whether the model matters or only
%     the coverage reward. This, not uniform random, is the honest baseline --- a
%     random baseline against a deep combination lock is a strawman, and an
%     AFL-style climber walks such a lock byte by byte,
%   - an oracle policy: the ceiling.
%
% State which baseline each claim is made against, and never make a claim against the
% floor alone.

\section{Currencies: Executions versus Wall-Clock}\label{sec:methodology-currencies}

% TODO: Name the accounting problem explicitly and commit to a rule. MCTS spends many
% model evaluations per action plus training time, while a conventional fuzzer spends
% almost nothing per execution. Executions-to-solve flatters this approach;
% wall-clock flatters the fuzzer. Report BOTH for every headline claim and state which
% currency each claim rests on. Include the compute split (emulation, inference,
% training) from \cref{sec:platform-throughput}.
%
% Note where the corpus formulation changes this accounting: the operator sweep of
% \cref{sec:corpus-determinism} buys many executions per agent decision, so the two
% formulations sit at different points on the same trade-off.

\section{Multi-Run Protocol}\label{sec:methodology-runs}

% TODO: The reliability protocol, and the statistical framing it forces.
%
% Framing first: the outcome is bimodal, not Gaussian. A run either learns the task or
% collapses, so this is a mixture of two populations and a mean lands in a valley
% where no run lives. Report solve rate with a Wilson score interval, report
% conditional-on-success metrics separately ("of the runs that solved, the median
% executions to solve was ..."), and plot medians with quantile bands or all runs
% individually --- never mean and standard deviation.
%
% Because the environment is deterministic (\cref{sec:platform-determinism}), all
% run-to-run variance is algorithmic: initialization, search noise, replay sampling.
% This is a reliability study of the learner, and stating that makes the variance
% analysis of \cref{sec:results-ablations} interpretable rather than apologetic.
%
% Protocol: a pre-registered set of RNG seeds fixed before any results are looked at;
% at least ten runs for headline configurations, around five paired runs for
% load-bearing ablations, three for peripheral ones and labelled indicative. Compare
% paired at fixed RNG seeds --- McNemar or a paired bootstrap for rates, Wilcoxon signed
% rank for continuous metrics --- and report effect sizes rather than p-values alone.
% If an effect falls inside the run-to-run band, report insensitivity: that is a result.
%
% Distinguish "ever reached" from "still solving at the end of training": late
% collapse wipes out earlier success, so state when the reported model was snapshotted.

\section{Hyperparameter Justification}\label{sec:methodology-hyperparameters}

% TODO: The question to answer is not "is this optimal" but "are the conclusions an
% artifact of this choice". Tier every hyperparameter and justify it accordingly:
%   - inherited or incidental: cite the source, one sentence,
%   - adapted by necessity: derived from problem structure. The value support is the
%     worked example --- the reward gap must span more than one bin or the heads go
%     action-blind (\cref{sec:input-reward}); show the arithmetic,
%   - sensitivity-analyzed: the few parameters the result hinges on, with a small
%     multi-run sweep showing a stable basin.
%
% State that no joint grid search was run and why (runs cost hours; run-to-run variance can
% swamp the parameter effect, so single-seed sweeps mislead). For the learning rate,
% use the evidence already in hand rather than inheritance: a rate that destabilized
% training through invalid replay priorities is a stronger justification than a
% citation. Any reported sweep uses at least three runs per point.
%
% Finally, the move that answers "you tuned to your result": the configuration is
% frozen after development on \cref{sec:target-sequence} and transferred UNCHANGED to
% every other target. Report the frozen configuration once, as a table, and say when
% it was frozen. There is no train/test split over environments; cross-target
% transfer is the substitute.

\section{Exploratory and Confirmatory Results}\label{sec:methodology-preregistration}

% TODO: Separate the two, explicitly. The development history was exploratory and
% hypothesis-generating; anything discovered by that search is contaminated as a
% tested claim, and with dozens of diagnostic series some will look significant by
% chance. Every headline claim in \cref{chap:results} is re-run fresh, on the
% pre-registered set of RNG seeds, with the frozen configuration and the final code, and
% reported as confirmatory. Results from before a correctness fix are never mixed with
% results from after it.
